\section{Fibre decoding and target observability}
\label{sec:fibre-decoding}

Fix an anchor assignment.  The remaining CRT coordinates are indexed by
\[
 R_{\mathrm{rows}}=(P\setminus B)\sqcup S_b.
\]
For $r\in R_{\mathrm{rows}}$, a residue $x\pmod r$, and $q\in B$, let $H_{rq}(x,a_q)$ be the centred CRT lift of $x\pmod r$ and $a_q\pmod q$, and put
\[
 u_{rq}(x)=\frac{H_{rq}(x,a_q)}{rq}\pmod1.
\]
If the row residue changes by $d$, then
\begin{equation}
\label{eq:row-shift}
 u_{rq}(x+d)-u_{rq}(x)=dq^{-1}/r\pmod1.
\end{equation}
Thus the separation of two possible row residues is independent of the anchor shift.

\subsection{Cyclic row distance}

\begin{lemma}[Multiplicity-sensitive cyclic energy]
\label{lem:cyclic-energy}
Let $c_1,\ldots,c_M$ be nonzero residues modulo a prime $r$, each occurring at most $\mu$ times.  If $M/\mu$ is sufficiently large, then for every $d\not\equiv0\pmod r$,
\[
 \sum_{j=1}^M\dist{dc_j/r}^2
 \geq c\frac{M^3}{\mu^2r^2}
\]
for an absolute constant $c>0$.
\end{lemma}

\begin{proof}
Multiplication by the nonzero residue $d$ permutes the nonzero residue classes modulo $r$ and therefore preserves the multiplicity bound.  Put
\[
 u=\left\lfloor\frac{M}{4\mu}\right\rfloor.
\]
For sufficiently large $M/\mu$, one has $u\geq M/(8\mu)\geq1$.  At most $2\mu u\leq M/2$ entries can lie in the $2u$ nonzero residue classes closest to zero on the circle.  Every remaining entry has circle distance at least
\[
 \frac{u}{r}\geq\frac{M}{8\mu r}.
\]
At least $M/2$ entries therefore contribute the square of this quantity, which proves the stated lower bound.
\end{proof}

For a row $r$, define
\[
 D_r=\min_{d\not\equiv0\,(r)}\sum_{q\in B}\dist{dq^{-1}/r}^2.
\]

\begin{proposition}[Lower-prime row distance]
\label{prop:lower-row-distance}
Uniformly for every prime $r\in P\setminus B$,
\[
 D_r\gg\frac{Z}{\log^3Z}.
\]
\end{proposition}

\begin{proof}
Inversion permutes the nonzero residue classes modulo $r$.  The interval $[Z/2,Z)$ contains at most
\[
 \mu_r\leq Z/(2r)+1\leq2Z/r
\]
integers in one residue class modulo $r$.  With $M=|B|\gg Z/\log Z$, \cref{lem:cyclic-energy} gives
\[
 D_r\gg\frac{M^3}{\mu_r^2r^2}
 \gg\frac{Z^3/\log^3Z}{Z^2}
 =\frac{Z}{\log^3Z}.
\]
Moreover $M/\mu_r\gg r/\log Z\geq X/\log Z$, so the largeness hypothesis is uniform in $r$.
\end{proof}

\begin{proposition}[Target-row distance]
\label{prop:target-row-distance}
For every $r\in S_b$,
\[
 D_r\geq\frac{|B|}{r^2}\gg_b\frac{Z}{\log Z}.
\]
\end{proposition}

\begin{proof}
Every $q\in B$ is coprime to $r$.  If $d\not\equiv0\pmod r$, then $dq^{-1}$ is nonzero modulo $r$, so each circle distance is at least $1/r$.
\end{proof}

The target-row sensors depend only on the reduced denominator $b$.  Their common kernel on $\prod_{r\in S_b}\mathbb Z/r\mathbb Z$ is trivial.  The numerator $a$ of the fixed target does not enter these sensors and creates no additional decoder or CRT obstruction.

\subsection{Row tails}

For $r\in R_{\mathrm{rows}}$, define
\[
 E_r(x)=\sum_{q\in B}\dist{u_{rq}(x)}^2,
 \qquad
 A_r(x)=\prod_{q\in B}
 \left|(1-\theta)+\theta\exp(2\pi i u_{rq}(x))\right|,
\]
and choose $x_r^*$ minimizing $E_r(x)$.

\begin{lemma}[Shift-uniform row tail]
\label{lem:row-tail}
There is
\[
 c_{\rm row}=c_{\rm row}(t,\gamma,b)>0
\]
such that
\[
 \delta_r:=\sum_{x\neq x_r^*}A_r(x)
 \leq r\exp(-c_{\rm row}D_r).
\]
\end{lemma}

\begin{proof}
If $x\neq x_r^*$, \cref{eq:row-shift} and the circle triangle inequality imply
\[
 D_r\leq2E_r(x)+2E_r(x_r^*)\leq4E_r(x).
\]
Apply \cref{eq:bernoulli-modulus} with $\kappa_{\rm mod}$ and sum over the at most $r$ nondecoder residues.  One may take $c_{\rm row}=\kappa_{\rm mod}/4$.
\end{proof}

Consequently,
\begin{equation}
\label{eq:Delta-row}
 \Delta:=\sum_{r\in R_{\mathrm{rows}}}\delta_r
 \leq Z^2\exp(-c_{\rm row}Z/\log^3Z)=o(1).
\end{equation}

\subsection{The exact retained skeleton}

For the fixed anchor assignment, assign every factor involving one anchor prime and one row prime to the corresponding row modulus $A_r$.  Internal anchor factors occur in $T_B$, while every retained lower--lower factor and every phase remains in a residual factor.  Thus
\begin{equation}
\label{eq:concrete-fibre-factorization}
 F(a_B,x_R)
 =T_B(a_B)\prod_{r\in R_{\mathrm{rows}}}A_r(x_r)
  H_{a_B}(x_R),
 \qquad |H_{a_B}(x_R)|\leq1.
\end{equation}
If the full modulus vanishes, set $H_{a_B}=0$.  The exact weighted compression theorem and \cref{eq:P-top,eq:Delta-row} give:

\begin{proposition}[Global weighted fibre error]
\label{prop:global-fibre-error}
The total contribution of every nondecoder fibre point is at most
\begin{equation}
\label{eq:global-fibre-error}
 P_{\mathrm{anc}}\bigl(\exp(\Delta)-1\bigr)
 \ll Z\log Z\,\Delta=o(1).
\end{equation}
\end{proposition}

No decoder weight is divided out, and the raw number of anchor assignments does not occur.

\subsection{Identification ranges}

Suppose the anchor assignment is coherent with integer label $m$.  The candidate row residue $m\pmod r$ satisfies
\begin{equation}
\label{eq:candidate-row-energy}
 E_r(m\bmod r)
 \leq\frac{m^2}{r^2}W_B
 \ll\frac{m^2/r^2}{Z\log Z}.
\end{equation}
If this is less than $D_r/8$, the candidate is the unique minimizer.

Put
\begin{equation}
\label{eq:M-dec}
 M_{\mathrm{dec}}=\frac{XZ}{(\log Z)^2}.
\end{equation}

\begin{proposition}[Prime-coordinate identification]
\label{prop:prime-decoder-identification}
For every $r\in P\setminus B$ and every coherent anchor label satisfying $|m|\leq M_{\mathrm{dec}}$,
\[
 x_r^*=m\pmod r.
\]
\end{proposition}

\begin{proof}
At $r=X$ and $|m|=M_{\mathrm{dec}}$, \cref{eq:candidate-row-energy} is $O(Z/\log^5Z)$, whereas \cref{prop:lower-row-distance} gives $D_r\gg Z/\log^3Z$.
\end{proof}

Fix once and for all a sufficiently small constant
\[
 \kappa_0=\kappa_0(t,\gamma,b)>0,
 \qquad \kappa_0<\frac14,
\]
small enough for the finitely many target rows and for the common linear phase range below.  Define the universal full-coordinate cutoff
\begin{equation}
\label{eq:T-full}
 T_0=\kappa_0\min(X^2,Z).
\end{equation}

\begin{proposition}[Target-coordinate identification]
\label{prop:target-decoder-identification}
For every $r\in S_b$ and every coherent anchor label satisfying $|m|\leq T_0$,
\[
 x_r^*=m\pmod r.
\]
\end{proposition}

\begin{proof}
By \cref{eq:candidate-row-energy}, the candidate energy is $O_b(m^2/(Z\log Z))$, while \cref{prop:target-row-distance} is $\gg_b Z/\log Z$.  Since $|m|\leq T_0\leq\kappa_0Z$, choosing $\kappa_0$ sufficiently small makes the comparison strict, uniformly over the finite set $S_b$.
\end{proof}

Because $b$ is squarefree, agreement modulo every $r\in S_b$ is agreement modulo $b$.  Hence for $|m|\leq T_0$, the entire decoded tuple is the genuine integer $m$ modulo $L$.  In the larger range $T_0<|m|\leq M_{\mathrm{dec}}$, all $P$-coordinates still decode to $m$, while no target-coordinate identification is asserted.  The transition and adaptive estimates use only those prime coordinates.
